\chapter*{Resumen}
\addcontentsline{toc}{chapter}{Resumen}

SmartSense Monitoring es una plataforma abierta, de operación local, para registrar la temperatura superficial de equipos industriales y apoyar su monitoreo de condición. Tres nodos inalámbricos, cada uno con una sonda PT100 de tres hilos, un acondicionador MAX31865, un microcontrolador ESP32-C6 y un radio E220, envían sus lecturas a un gateway que conserva el historial y lo presenta en un dashboard local, sin servicio en la nube ni conexión permanente a Internet. Durante el piloto se reunieron \TotalMuestras{} muestras: los nodos A y C cubren \DuracionComunDias{} días, del 8 de mayo al 9 de junio de 2026, y el nodo B va del 10 de mayo al 20 de julio; las comparaciones entre nodos se hacen sobre el tramo en que coinciden. La disponibilidad cruda de datos fue de 76.0, 32.3 y 75.1\,\% por nodo, y de \DisponibilidadEnlace{} en los periodos clasificados como activos. El PDR estimado a partir de los intervalos de llegada fue de \PDREstimado{}. Cinco parejas de lecturas frente a un termómetro infrarrojo muestran concordancia de campo tras el ajuste de offset; la exactitud absoluta queda pendiente de una referencia calibrada. Una descarga parcial de cinco días permite estimar la autonomía en \AutonomiaLinealEstimada{} con un modelo lineal y en \AutonomiaModeloEstimada{} con la curva de descarga LiPo; ambas son estimaciones, no una duración medida. El sistema apoya el monitoreo de condición y no implementa un algoritmo de mantenimiento predictivo.

\noindent\textbf{Palabras clave:} monitoreo de condición, temperatura superficial, PT100, LoRa, plataforma local, Internet industrial de las cosas.

\clearpage
\chapter*{Abstract}
\addcontentsline{toc}{chapter}{Abstract}

SmartSense Monitoring is an open, locally operated platform that records the surface temperature of industrial equipment to support condition monitoring. Three wireless nodes, each with a three-wire PT100 probe, a MAX31865 front end, an ESP32-C6 microcontroller and an E220 radio, send their readings to a gateway that keeps the history and shows it on a local dashboard, with no cloud service or permanent Internet connection. The pilot collected \TotalMuestras{} samples: nodes A and C cover \DuracionComunDias{} days, from 8 May to 9 June 2026, and node B runs from 10 May to 20 July; comparisons between nodes use the interval they share. Raw data availability was 76.0, 32.3 and 75.1\,\% per node, and \DisponibilidadEnlace{} over the periods classified as active. The PDR estimated from arrival intervals was \PDREstimado{}. Five paired readings against an infrared thermometer show field agreement after offset adjustment; absolute accuracy awaits a calibrated reference. A five-day partial discharge gives an estimated autonomy of 177 days with a linear model and 291--307 days with a LiPo discharge curve; both are estimates, not a measured lifetime. The system supports condition monitoring and does not implement a predictive-maintenance algorithm.

\noindent\textbf{Keywords:} condition monitoring, surface temperature, PT100, LoRa, local platform, industrial Internet of Things.
\clearpage
